\IfFileExists{stacks-project.cls}{%
\documentclass{stacks-project}
}{%
\documentclass{amsart}
}

% For dealing with references we use the comment environment
\usepackage{verbatim}
\newenvironment{reference}{\comment}{\endcomment}
%\newenvironment{reference}{}{}
\newenvironment{slogan}{\comment}{\endcomment}
\newenvironment{history}{\comment}{\endcomment}

% For commutative diagrams we use Xy-pic
\usepackage[all]{xy}

% We use 2cell for 2-commutative diagrams.
\xyoption{2cell}
\UseAllTwocells

% We use multicol for the list of chapters between chapters
\usepackage{multicol}

% This is generally recommended for better output
\usepackage{lmodern}
\usepackage[T1]{fontenc}

% For cross-file-references
\usepackage{xr-hyper}

% Package for hypertext links:
\usepackage{hyperref}

% For any local file, say "hello.tex" you want to link to please
% use \externaldocument[hello-]{hello}
\externaldocument[introduction-]{introduction}
\externaldocument[conventions-]{conventions}
\externaldocument[sets-]{sets}
\externaldocument[categories-]{categories}
\externaldocument[topology-]{topology}
\externaldocument[sheaves-]{sheaves}
\externaldocument[sites-]{sites}
\externaldocument[stacks-]{stacks}
\externaldocument[fields-]{fields}
\externaldocument[algebra-]{algebra}
\externaldocument[brauer-]{brauer}
\externaldocument[homology-]{homology}
\externaldocument[derived-]{derived}
\externaldocument[simplicial-]{simplicial}
\externaldocument[more-algebra-]{more-algebra}
\externaldocument[smoothing-]{smoothing}
\externaldocument[modules-]{modules}
\externaldocument[sites-modules-]{sites-modules}
\externaldocument[injectives-]{injectives}
\externaldocument[cohomology-]{cohomology}
\externaldocument[sites-cohomology-]{sites-cohomology}
\externaldocument[dga-]{dga}
\externaldocument[dpa-]{dpa}
\externaldocument[sdga-]{sdga}
\externaldocument[hypercovering-]{hypercovering}
\externaldocument[schemes-]{schemes}
\externaldocument[constructions-]{constructions}
\externaldocument[properties-]{properties}
\externaldocument[morphisms-]{morphisms}
\externaldocument[coherent-]{coherent}
\externaldocument[divisors-]{divisors}
\externaldocument[limits-]{limits}
\externaldocument[varieties-]{varieties}
\externaldocument[topologies-]{topologies}
\externaldocument[descent-]{descent}
\externaldocument[perfect-]{perfect}
\externaldocument[more-morphisms-]{more-morphisms}
\externaldocument[flat-]{flat}
\externaldocument[groupoids-]{groupoids}
\externaldocument[more-groupoids-]{more-groupoids}
\externaldocument[etale-]{etale}
\externaldocument[chow-]{chow}
\externaldocument[intersection-]{intersection}
\externaldocument[pic-]{pic}
\externaldocument[weil-]{weil}
\externaldocument[adequate-]{adequate}
\externaldocument[dualizing-]{dualizing}
\externaldocument[duality-]{duality}
\externaldocument[discriminant-]{discriminant}
\externaldocument[derham-]{derham}
\externaldocument[local-cohomology-]{local-cohomology}
\externaldocument[algebraization-]{algebraization}
\externaldocument[curves-]{curves}
\externaldocument[resolve-]{resolve}
\externaldocument[models-]{models}
\externaldocument[functors-]{functors}
\externaldocument[equiv-]{equiv}
\externaldocument[pione-]{pione}
\externaldocument[etale-cohomology-]{etale-cohomology}
\externaldocument[proetale-]{proetale}
\externaldocument[relative-cycles-]{relative-cycles}
\externaldocument[more-etale-]{more-etale}
\externaldocument[trace-]{trace}
\externaldocument[crystalline-]{crystalline}
\externaldocument[spaces-]{spaces}
\externaldocument[spaces-properties-]{spaces-properties}
\externaldocument[spaces-morphisms-]{spaces-morphisms}
\externaldocument[decent-spaces-]{decent-spaces}
\externaldocument[spaces-cohomology-]{spaces-cohomology}
\externaldocument[spaces-limits-]{spaces-limits}
\externaldocument[spaces-divisors-]{spaces-divisors}
\externaldocument[spaces-over-fields-]{spaces-over-fields}
\externaldocument[spaces-topologies-]{spaces-topologies}
\externaldocument[spaces-descent-]{spaces-descent}
\externaldocument[spaces-perfect-]{spaces-perfect}
\externaldocument[spaces-more-morphisms-]{spaces-more-morphisms}
\externaldocument[spaces-flat-]{spaces-flat}
\externaldocument[spaces-groupoids-]{spaces-groupoids}
\externaldocument[spaces-more-groupoids-]{spaces-more-groupoids}
\externaldocument[bootstrap-]{bootstrap}
\externaldocument[spaces-pushouts-]{spaces-pushouts}
\externaldocument[spaces-chow-]{spaces-chow}
\externaldocument[groupoids-quotients-]{groupoids-quotients}
\externaldocument[spaces-more-cohomology-]{spaces-more-cohomology}
\externaldocument[spaces-simplicial-]{spaces-simplicial}
\externaldocument[spaces-duality-]{spaces-duality}
\externaldocument[formal-spaces-]{formal-spaces}
\externaldocument[restricted-]{restricted}
\externaldocument[spaces-resolve-]{spaces-resolve}
\externaldocument[formal-defos-]{formal-defos}
\externaldocument[defos-]{defos}
\externaldocument[cotangent-]{cotangent}
\externaldocument[examples-defos-]{examples-defos}
\externaldocument[algebraic-]{algebraic}
\externaldocument[examples-stacks-]{examples-stacks}
\externaldocument[stacks-sheaves-]{stacks-sheaves}
\externaldocument[criteria-]{criteria}
\externaldocument[artin-]{artin}
\externaldocument[quot-]{quot}
\externaldocument[stacks-properties-]{stacks-properties}
\externaldocument[stacks-morphisms-]{stacks-morphisms}
\externaldocument[stacks-limits-]{stacks-limits}
\externaldocument[stacks-cohomology-]{stacks-cohomology}
\externaldocument[stacks-perfect-]{stacks-perfect}
\externaldocument[stacks-introduction-]{stacks-introduction}
\externaldocument[stacks-more-morphisms-]{stacks-more-morphisms}
\externaldocument[stacks-geometry-]{stacks-geometry}
\externaldocument[moduli-]{moduli}
\externaldocument[moduli-curves-]{moduli-curves}
\externaldocument[examples-]{examples}
\externaldocument[exercises-]{exercises}
\externaldocument[guide-]{guide}
\externaldocument[desirables-]{desirables}
\externaldocument[coding-]{coding}
\externaldocument[obsolete-]{obsolete}
\externaldocument[fdl-]{fdl}
\externaldocument[index-]{index}

% Theorem environments.
%
\theoremstyle{plain}
\newtheorem{theorem}[subsection]{Theorem}
\newtheorem{proposition}[subsection]{Proposition}
\newtheorem{lemma}[subsection]{Lemma}

\theoremstyle{definition}
\newtheorem{definition}[subsection]{Definition}
\newtheorem{example}[subsection]{Example}
\newtheorem{exercise}[subsection]{Exercise}
\newtheorem{situation}[subsection]{Situation}

\theoremstyle{remark}
\newtheorem{remark}[subsection]{Remark}
\newtheorem{remarks}[subsection]{Remarks}

\numberwithin{equation}{subsection}

% Macros
%
\def\lim{\mathop{\mathrm{lim}}\nolimits}
\def\colim{\mathop{\mathrm{colim}}\nolimits}
\def\Spec{\mathop{\mathrm{Spec}}}
\def\Hom{\mathop{\mathrm{Hom}}\nolimits}
\def\Ext{\mathop{\mathrm{Ext}}\nolimits}
\def\SheafHom{\mathop{\mathcal{H}\!\mathit{om}}\nolimits}
\def\SheafExt{\mathop{\mathcal{E}\!\mathit{xt}}\nolimits}
\def\Sch{\mathit{Sch}}
\def\Mor{\mathop{\mathrm{Mor}}\nolimits}
\def\Ob{\mathop{\mathrm{Ob}}\nolimits}
\def\Sh{\mathop{\mathit{Sh}}\nolimits}
\def\NL{\mathop{N\!L}\nolimits}
\def\CH{\mathop{\mathrm{CH}}\nolimits}
\def\proetale{{pro\text{-}\acute{e}tale}}
\def\etale{{\acute{e}tale}}
\def\QCoh{\mathit{QCoh}}
\def\Ker{\mathop{\mathrm{Ker}}}
\def\Im{\mathop{\mathrm{Im}}}
\def\Coker{\mathop{\mathrm{Coker}}}
\def\Coim{\mathop{\mathrm{Coim}}}

% Boxtimes
%
\DeclareMathSymbol{\boxtimes}{\mathbin}{AMSa}{"02}

%
% Macros for moduli stacks/spaces
%
\def\QCohstack{\mathcal{QC}\!\mathit{oh}}
\def\Cohstack{\mathcal{C}\!\mathit{oh}}
\def\Spacesstack{\mathcal{S}\!\mathit{paces}}
\def\Quotfunctor{\mathrm{Quot}}
\def\Hilbfunctor{\mathrm{Hilb}}
\def\Curvesstack{\mathcal{C}\!\mathit{urves}}
\def\Polarizedstack{\mathcal{P}\!\mathit{olarized}}
\def\Complexesstack{\mathcal{C}\!\mathit{omplexes}}
% \Pic is the operator that assigns to X its picard group, usage \Pic(X)
% \Picardstack_{X/B} denotes the Picard stack of X over B
% \Picardfunctor_{X/B} denotes the Picard functor of X over B
\def\Pic{\mathop{\mathrm{Pic}}\nolimits}
\def\Picardstack{\mathcal{P}\!\mathit{ic}}
\def\Picardfunctor{\mathrm{Pic}}
\def\Deformationcategory{\mathcal{D}\!\mathit{ef}}
\usepackage{cleveref}
\crefname{originalequation}{original equation}{original equations}
\Crefname{originalequation}{Original equation}{Original equations}
\makeatletter
\newlabel{equation-extended-alternating}{{0G7J}{source}{}{sourcepointer-1}{}}
\newlabel{equation-extended-alternating@cref}{{[originalequation][0][]0G7J}{[1][0][]source}{}{}{}}
\makeatother
\crefname{originallemma}{original lemma}{original lemmas}
\Crefname{originallemma}{Original lemma}{Original lemmas}
\makeatletter
\newlabel{lemma-extended-alternating-represented}{{0G7M}{source}{}{sourcepointer-2}{}}
\newlabel{lemma-extended-alternating-represented@cref}{{[originallemma][0][]0G7M}{[1][0][]source}{}{}{}}
\makeatother
\crefname{originalremark}{original remark}{original remarks}
\Crefname{originalremark}{Original remark}{Original remarks}
\makeatletter
\newlabel{remark-support-c-equations}{{0G7I}{source}{}{sourcepointer-3}{}}
\newlabel{remark-support-c-equations@cref}{{[originalremark][0][]0G7I}{[1][0][]source}{}{}{}}
\makeatother
\crefname{originalremark}{original remark}{original remarks}
\Crefname{originalremark}{Original remark}{Original remarks}
\makeatletter
\newlabel{remark-supported-map-c-equations}{{0G7Q}{source}{}{sourcepointer-4}{}}
\newlabel{remark-supported-map-c-equations@cref}{{[originalremark][0][]0G7Q}{[1][0][]source}{}{}{}}
\makeatother
\title{Determinant transition formula for the local cohomology comparison maps under a change of generators}
\author{The Stacks Project Authors}
\date{stacks-project:perfect.tex:local-snapshot-20261012; GNU FDL 1.2 or later}
\begin{document}
\maketitle
\noindent\textit{Editorial note.} The source ends 'Some details omitted' inside the Koszul-differential verification; the omitted step is the term-by-term cancellation argument for the alternating sum, carried as an explicit obligation.\par\medskip
\noindent\textit{Editorial note.} Inputs lemma-extended-alternating-represented and lemma-application-nice-K-injective are same-chapter NOT carded (cross-line-claimed labels) and are declared external.\par\medskip
\noindent\textit{Editorial note.} Steps the author leaves to the reader are flagged here for the reviewer: source line(s) 1442 carry an "omitted"/"follows easily"/"immediate from the definitions" phrase. Those sentences are retained verbatim; the displayed argument around them stays inside the counted span.\par\medskip
\noindent\textit{Allowed standard prerequisites (editorial).}\begin{itemize}
\item Extended alternating complexes represent local cohomology. lemma-extended-alternating-represented (same chapter, NOT carded, declared external). Identifies the three Cech complexes with i\_*R H\_Z(F) up to canonical maps representing the identity.
\item Local comparison maps from generators. Remark remark-supported-map-c-equations (same chapter, structural). Defines the maps c\_\{f\_1,...,f\_c\} attached to a regular sequence cutting out Z.
\item Functoriality of Cech complexes under the subdivision. lemma-application-nice-K-injective (same chapter, NOT carded, declared external). Underlies the compatibility used when comparing the three complexes.
\end{itemize}\medskip
\ifdefined\hypertarget\hypertarget{sourcepointer-1}{}\fi
\noindent\textit{Source reference.} Original equation 0G7J: sources/stacks-perfect/pinned-original-source.tex, source lines 1107--1107. Navigation pointer only: the counted proof cites this statement by \textbackslash{}ref as a standing input; the deduction the counted proof has to make is carried out inside the counted span, and no part of that deduction is delegated to this pointer.\par\medskip
\ifdefined\hypertarget\hypertarget{sourcepointer-2}{}\fi
\noindent\textit{Source reference.} Original lemma 0G7M: sources/stacks-perfect/pinned-original-source.tex, source lines 1185--1185. Navigation pointer only: the counted proof cites this statement by \textbackslash{}ref as a standing input; the deduction the counted proof has to make is carried out inside the counted span, and no part of that deduction is delegated to this pointer.\par\medskip
\ifdefined\hypertarget\hypertarget{sourcepointer-3}{}\fi
\noindent\textit{Source reference.} Original remark 0G7I: sources/stacks-perfect/pinned-original-source.tex, source lines 1093--1093. Navigation pointer only: the counted proof cites this statement by \textbackslash{}ref as a standing input; the deduction the counted proof has to make is carried out inside the counted span, and no part of that deduction is delegated to this pointer.\par\medskip
\ifdefined\hypertarget\hypertarget{sourcepointer-4}{}\fi
\noindent\textit{Source reference.} Original remark 0G7Q: sources/stacks-perfect/pinned-original-source.tex, source lines 1267--1267. Navigation pointer only: the counted proof cites this statement by \textbackslash{}ref as a standing input; the deduction the counted proof has to make is carried out inside the counted span, and no part of that deduction is delegated to this pointer.\par\medskip
\begin{lemma}
\label{lemma-supported-map-determinant}
With $X$, $f_1, \ldots, f_c \in \Gamma(X, \mathcal{O}_X)$, and
$\mathcal{F}$ as in Remark \ref{remark-support-c-equations}.
Let $a_{ji} \in \Gamma(X, \mathcal{O}_X)$ for $1 \leq i, j \leq c$
and set $g_j = \sum_{i = 1, \ldots, c} a_{ji}f_i$. Assume $g_1, \ldots, g_c$
scheme theoretically cut out $Z$. If $\mathcal{F}$ is quasi-coherent, then
$$
c_{f_1, \ldots, f_c} = \det(a_{ji}) c_{g_1, \ldots, g_c}
$$
where $c_{f_1, \ldots, f_c}$ and $c_{g_1, \ldots, g_c}$ are as in
Remark \ref{remark-supported-map-c-equations}.
\end{lemma}

\begin{proof}
We will prove that $c_{f_1, \ldots, f_c}(s) =
\det(a_{ij}) c_{g_1, \ldots, g_c}(s)$ as global sections of
$\mathcal{H}_Z(\mathcal{F})$ for any $s \in \mathcal{F}(X)$.
This is sufficient since we then obtain the same result for
section over any open subscheme of $X$.
To do this, for $1 \leq i_0 < \ldots < i_p \leq c$
and $1 \leq j_0 < \ldots < j_q \leq c$ we denote
$U_{i_0 \ldots i_p} \subset X$,
$V_{j_0 \ldots j_q} \subset X$, and
$W_{i_0 \ldots i_p, j_0 \ldots j_q} \subset X$
the open subscheme where $f_{i_0} \ldots f_{i_p}$ is invertible,
$g_{j_0} \ldots g_{j_q}$ is invertible, and where
$f_{i_0} \ldots f_{i_p}g_{j_0} \ldots g_{j_q}$ is invertible.
We denote $\mathcal{F}_{i_0 \ldots i_p}$,
resp.\ $\mathcal{F}'_{j_0 \ldots j_q}$
$\mathcal{F}''_{i_0 \ldots i_p, j_0 \ldots j_q}$
the pushforward to $X$ of the restriction of $\mathcal{F}$ to
$U_{i_0 \ldots i_p}$, resp.\ $V_{j_0 \ldots j_q}$,
resp.\ $W_{i_0 \ldots i_p, j_0 \ldots j_q}$.
Then we obtain three extended alternating {\v C}ech complexes
$$
\mathcal{F}^\bullet :
\mathcal{F} \to \bigoplus\nolimits_{i_0} \mathcal{F}_{i_0} \to
\bigoplus\nolimits_{i_0 < i_1} \mathcal{F}_{i_0i_1} \to \ldots
$$
and
$$
(\mathcal{F}')^\bullet :
\mathcal{F} \to \bigoplus\nolimits_{j_0} \mathcal{F}'_{j_0} \to
\bigoplus\nolimits_{j_0 < j_1} \mathcal{F}'_{j_0j_1} \to \ldots
$$
and
$$
(\mathcal{F}'')^\bullet :
\mathcal{F} \to
\bigoplus\nolimits_{i_0} \mathcal{F}_{i_0} \oplus
\bigoplus\nolimits_{j_0} \mathcal{F}'_{j_0} \to
\bigoplus\nolimits_{i_0 < i_1} \mathcal{F}_{i_0i_1} \oplus
\bigoplus\nolimits_{i_0, j_0} \mathcal{F}''_{i_0, j_0} \oplus
\bigoplus\nolimits_{j_0 < j_1} \mathcal{F}'_{j_0j_1} \to
\ldots
$$
whose differentials are those used in defining
(\ref{equation-extended-alternating}).
There are maps of complexes
$$
(\mathcal{F}'')^\bullet \to \mathcal{F}^\bullet
\quad\text{and}\quad
(\mathcal{F}'')^\bullet \to (\mathcal{F}')^\bullet
$$
given by the projection maps on the terms (and hence inducing
the identity map in degree $0$). Observe that by
Lemma \ref{lemma-extended-alternating-represented}
each of these complexes represents
$i_*R\mathcal{H}_Z(\mathcal{F})$ and these maps
represent the identity on this object. Thus it suffices
to find an element
$$
\sigma \in H^c((\mathcal{F}'')^\bullet(X))
$$
mapping to $c_{f_1, \ldots, f_c}(s)$ and $\det(a_{ji})c_{g_1, \ldots, g_c}(s)$
by these two maps. It turns out we can explicitly give a cocycle
for $\sigma$. Namely, we take
$$
\sigma_{1 \ldots c} = \frac{s}{f_1 \ldots f_c} \in \mathcal{F}_{1 \ldots c}(X)
\quad\text{and}\quad
\sigma'_{1 \ldots c} = \frac{\det(a_{ji})s}{g_1 \ldots g_c} \in
\mathcal{F}'_{1 \ldots c}(X)
$$
and we take
$$
\sigma_{i_0 \ldots i_p, j_0 \ldots j_{c - p - 2}} =
\frac{\lambda(i_0 \ldots i_p, j_0 \ldots j_{c - p - 2})s}%
{f_{i_0} \ldots f_{i_p}g_{j_0} \ldots g_{j_{c - p - 2}}}
\in \mathcal{F}''_{i_0 \ldots i_p, j_0 \ldots j_{c - p - 2}}(X)
$$
where $\lambda(i_0 \ldots i_p, j_0 \ldots j_{c - p - 2})$
is the coefficient of $e_1 \wedge \ldots \wedge e_c$ in
the formal expression
$$
e_{i_0} \wedge \ldots \wedge e_{i_p} \wedge
(a_{j_01} e_1 + \ldots + a_{j_0c}e_c) \wedge \ldots \wedge
(a_{j_{c - p - 2}1} e_1 + \ldots + a_{j_{c - p - 2}c}e_c)
$$
To verify that $\sigma$ is a cocycle, we have to show for
$1 \leq i_0 < \ldots < i_p \leq c$ and
$1 \leq j_0 < \ldots < j_{c - p - 1} \leq c$
that we have
\begin{align*}
0 & =
\sum\nolimits_{a = 0, \ldots, p} (-1)^a
f_{i_a} \lambda(i_0 \ldots \hat i_a \ldots i_p, j_0 \ldots j_{c - p - 1}) \\
& +
\sum\nolimits_{b = 0, \ldots, c - p - 1} (-1)^{p + b + 1}g_{j_b}
\lambda(i_0 \ldots i_p, j_0 \ldots \hat j_b \ldots j_{c - p - 1})
\end{align*}
The easiest way to see this is perhaps to argue that the formal expression
$$
\xi = e_{i_0} \wedge \ldots \wedge e_{i_p} \wedge
(a_{j_01} e_1 + \ldots + a_{j_0c}e_c) \wedge \ldots \wedge
(a_{j_{c - p - 1}1} e_1 + \ldots + a_{j_{c - p - 1}c}e_c)
$$
is $0$ as it is an element of the $(c + 1)$st wedge power of the free module
on $e_1, \ldots, e_c$ and that the expression above is the image of
$\xi$ under the Koszul differential sending $e_i \to f_i$. Some details
omitted.
\end{proof}

\end{document}
